\documentclass[11pt]{article}

\usepackage[margin=1in]{geometry}
\usepackage{amsmath,amssymb,amsthm}
\usepackage{booktabs}
\usepackage{enumitem}
\usepackage[hypertexnames=false]{hyperref}
\usepackage{listings}
\usepackage{xcolor}

\newtheorem{definition}{Definition}
\newtheorem{theorem}{Theorem}
\newtheorem{corollary}{Corollary}
\newtheorem{proposition}{Proposition}

\lstset{
  language=Python,
  basicstyle=\ttfamily\small,
  keywordstyle=\color{blue!70!black},
  commentstyle=\color{green!50!black},
  stringstyle=\color{red!60!black},
  showstringspaces=false,
  frame=single,
  breaklines=true,
  columns=fullflexible,
}

\title{\textbf{Convergence by Composition}\\
\large A Structured Adapter Architecture for Multi-Agent System Integration\\[0.5em]
\normalsize Preprint -- Feedback Welcome}
\author{Bogdan Banu\\\texttt{bogdan@banu.be}}
\date{\today}

\begin{document}
\maketitle

%% ============================================================================
%% ABSTRACT
%% ============================================================================
\begin{abstract}
Agent orchestration frameworks---Swarms, DeerFlow, AnimaWorks, Ralph, and
A-Evolve---make fundamentally different architectural choices (graph-based,
event-driven, hierarchical, evolutionary) yet share structural concerns:
error amplification across handoffs, coordination overhead from tool
proliferation, and the absence of convergence guarantees for adaptive loops.
We present a structured adapter architecture centred on a single
intermediate representation, \texttt{ExternalTopology}, that lets Operon's
structural analysis layer analyze \emph{any} external agent topology
without importing the target framework.  Five adapters parse framework-specific configurations into this
common representation; four epistemic theorems serve as a structural linter;
four TLA+ specifications verify safety invariants; and a co-design
fixed-point proof (following Zardini) establishes that the adaptive assembly
loop converges within $\varepsilon$.  The architecture is implemented in 14
convergence modules with 216 convergence-specific tests as part of a broader
suite of 1{,}404 tests and 98 worked examples.
\end{abstract}

%% ============================================================================
%% 1. INTRODUCTION
%% ============================================================================
\section{Introduction}

The past two years have seen a proliferation of agent orchestration
frameworks, each embodying a different theory of how autonomous agents
should coordinate.  Swarms~\cite{swarms} provides graph-based workflows
with sequential, concurrent, and hierarchical scheduling.
DeerFlow, built on LangGraph, offers sub-agent harnesses with sandboxed
execution and progressive skill loading.
AnimaWorks~\cite{animaworks} introduces developmental priming with
heartbeat-driven memory consolidation.  Ralph organizes agents as ``hats''
connected by event-driven transitions with backpressure constraints.
A-Evolve implements an evolutionary Solve--Observe--Evolve--Gate--Reload
loop that mutates workspace snapshots under a fitness gate.

These frameworks differ in execution model, memory architecture, and
deployment infrastructure, yet they converge on a shared set of structural
concerns.  Every multi-agent system must manage error amplification across
sequential handoffs, balance coordination overhead against parallel
speedup, and decide when an adaptive loop has stabilized.  As Evans,
Bratton, and Ag\"{u}era y Arcas argue, intelligence is ``fundamentally
plural, social, and relational''~\cite{evans2026plural}; the engineering
challenge is to build, verify, and optimize the structural infrastructure
for such systems.

The Agentic Computation Graph (ACG) survey by Yue et
al.~\cite{yue2026acg} provides a comprehensive taxonomy of workflow
optimization methods---from static template search to in-execution
editing---but does not address formal verification or compositional
convergence proofs.  Kim et al.~\cite{kim2025scaling} identify scaling
laws for agent systems whose functional form is consistent with the
epistemic bounds derived in the present work.

Our contribution is a \emph{structured adapter architecture} that decouples
structural analysis from runtime execution:

\begin{enumerate}[leftmargin=*,itemsep=2pt]
  \item \textbf{ExternalTopology} as a framework-agnostic intermediate
    representation: five adapters parse five frameworks into a single
    frozen dataclass; all downstream analysis is source-agnostic.
  \item \textbf{Epistemic theorems as structural linter}: error
    amplification, sequential penalty, parallel acceleration, and tool
    density bounds are applied uniformly to any topology.
  \item \textbf{TLA+ verified safety invariants}: four specifications
    cover template exchange, developmental gating, convergence detection,
    and evolution gating, verified by the TLC model checker with
    small-model parameters.
  \item \textbf{Co-design convergence}: each adapter is a design problem
    in Zardini's co-design theory~\cite{zardini2023codesign};
    series/parallel composition and feedback fixed-point iteration prove
    that scoring stabilizes.
\end{enumerate}

The rest of this paper describes the architecture (Section~\ref{sec:arch}),
formal results (Section~\ref{sec:formal}), implementation
(Section~\ref{sec:impl}), related work (Section~\ref{sec:related}),
discussion (Section~\ref{sec:discussion}), and conclusion
(Section~\ref{sec:conclusion}).

%% ============================================================================
%% 2. ARCHITECTURE
%% ============================================================================
\section{Architecture}\label{sec:arch}

The convergence architecture is a five-layer stack (Figure~\ref{fig:stack}).
Each layer is independently useful: the cognitive layer can run without an
orchestration layer; the structural layer can analyze topologies without
executing them.

\begin{figure}[ht]
\centering
\begin{verbatim}
  +-----------------------------------------------+
  |  AnimaWorks (cognitive layer)                  |
  |  Identity, memory consolidation, priming       |
  +-----------------------------------------------+
  |  AsyncThink (thinking layer)                   |
  |  Fork/Join policy, concurrency optimization    |
  +-----------------------------------------------+
  |  DeerFlow / Swarms / Ralph (orchestration)     |
  |  Sub-agent harness, sandbox, skill loading     |
  +-----------------------------------------------+
  |  A-Evolve (evolution layer)                    |
  |  Solve-Observe-Evolve-Gate-Reload loop         |
  +-----------------------------------------------+
  |  Operon (structural layer)                     |
  |  Topology advice, epistemic analysis,          |
  |  pattern scoring, watcher, co-design proofs    |
  +-----------------------------------------------+
\end{verbatim}
\caption{Five-layer convergence stack.  Operon sits at the bottom,
providing structural analysis to every layer above.  Arrows flow upward
(analysis results) and downward (topology descriptions).}
\label{fig:stack}
\end{figure}

\subsection{ExternalTopology as Framework-Agnostic Intermediate Representation}

The key abstraction is a single frozen dataclass that every adapter
produces and every analysis function consumes:

\begin{lstlisting}[caption={ExternalTopology dataclass (convergence/types.py).},label=lst:topo]
@dataclass(frozen=True)
class ExternalTopology:
    source: str                          # "swarms"|"deerflow"|...
    pattern_name: str
    agents: tuple[dict[str, Any], ...]
    edges: tuple[tuple[str, str], ...]
    metadata: dict[str, Any] = field(default_factory=dict)
\end{lstlisting}

The agent specs use \texttt{dict[str, Any]} to accommodate heterogeneous
framework configurations.  Python type hints annotate the top-level
fields; the nested dicts carry no schema enforcement at runtime.
The \texttt{@dataclass(frozen=True)} decorator prevents reassignment
of top-level fields but does not deep-freeze nested mutable objects.
The \texttt{source} field is a tag, not a type discriminant---all
downstream code treats every \texttt{ExternalTopology} identically.
Agents are plain dicts with at least \texttt{name} and \texttt{role}
keys; edges are directed pairs.  This representation is serializable as
JSON, enabling analysis results to be stored, transmitted, and compared
across framework boundaries.

\subsection{Analysis Pipeline}

The function \texttt{analyze\_external\_topology(topology)} takes any
\texttt{ExternalTopology} and returns an \texttt{AdapterResult}:

\begin{lstlisting}[caption={AdapterResult dataclass (convergence/types.py).},label=lst:result]
@dataclass(frozen=True)
class AdapterResult:
    topology_advice: TopologyAdvice
    suggested_template: PatternTemplate | None
    warnings: tuple[str, ...]
    risk_score: float                    # [0.0, 1.0]
\end{lstlisting}

The pipeline constructs a wiring diagram from agents and edges, applies the
four epistemic theorems, generates structural warnings, and optionally
produces a \texttt{PatternTemplate} for the pattern library.

\subsection{Design Principles}

Three principles govern the adapter architecture:

\begin{enumerate}[leftmargin=*,itemsep=2pt]
  \item \textbf{Serializable dicts, zero external imports.}  Every adapter
    operates on plain dicts parsed from configuration files or API
    responses.  No adapter imports Swarms, DeerFlow, AnimaWorks, Ralph, or
    A-Evolve.
  \item \textbf{One parser, one analyzer.}  Adding a new framework means
    writing one \texttt{parse\_*} function that produces
    \texttt{ExternalTopology}; the analysis code is unchanged.
  \item \textbf{Bidirectional exchange.}  DeerFlow's Markdown skill files
    convert to \texttt{PatternTemplate} and vice versa, enabling
    cross-framework template sharing.
\end{enumerate}

%% ============================================================================
%% 3. FORMAL RESULTS
%% ============================================================================
\section{Formal Results}\label{sec:formal}

\subsection{Epistemic Theorems as Structural Linter}\label{sec:epistemic}

Operon derives four structural bounds from epistemic topology---the
study of how knowledge and uncertainty propagate through agent
communication graphs.  These bounds are applied to every
\texttt{ExternalTopology} as a ``structural linter,'' flagging
architectures that are likely to amplify errors or incur excessive
coordination overhead.

\begin{definition}[Error Amplification Bound]
For a sequential chain of $n$ agents, each with per-stage error rate
$\epsilon$, the cumulative error probability is bounded by:
\[
  P_{\mathrm{error}}(n) \;\leq\; 1 - (1 - \epsilon)^n
\]
When $n\epsilon \ll 1$, the linear approximation $P_{\mathrm{error}} \approx n\epsilon$
applies.  The linter warns when the chain depth produces
$P_{\mathrm{error}} > \tau$ for a configurable tolerance $\tau$.
\end{definition}

\begin{definition}[Sequential Penalty]
For a topology with $n_{\mathrm{seq}}$ sequential handoffs and total
agent count $N$, the sequential penalty ratio is:
\[
  \rho_{\mathrm{seq}} \;=\; \frac{n_{\mathrm{seq}}}{N}
\]
Values of $\rho_{\mathrm{seq}} > 0.7$ indicate that the topology is
predominantly sequential and may benefit from parallelization.
\end{definition}

\begin{definition}[Parallel Acceleration]
For a topology with $k$ independent branches of depths $d_1, \ldots, d_k$,
the critical-path latency under parallel execution is:
\[
  L_{\mathrm{par}} \;=\; \max_{i \in [k]} d_i
\]
The speedup over sequential execution is $S = \sum_i d_i / L_{\mathrm{par}}$.
The linter reports the achievable speedup and flags topologies where
$S < 1.2$ (parallel structure exists but provides negligible benefit).
\end{definition}

\noindent
Parallel acceleration is computed and reported separately by the
epistemic analysis layer but is not consumed by the current topology
recommender or the composite risk score.  The risk score focuses on the
three primary failure modes (error amplification, sequential overhead,
tool density) plus a topology-mismatch flag.

\begin{definition}[Tool Density]
For a topology with $T$ distinct tools across $N$ agents, the tool density
is $\delta = T / N$.  High density ($\delta > 3$) suggests coordination
overhead: each agent must manage many tool interfaces.  Low density
($\delta < 0.5$) suggests under-utilization.
\end{definition}

The composite risk score is a weighted sum of four components:
\[
  r \;=\; 0.4\, \hat{P}_{\mathrm{error}}
    \;+\; 0.3\, \rho_{\mathrm{seq}}
    \;+\; 0.2\, \delta_{\mathrm{density}}
    \;+\; 0.1\, M_{\mathrm{topo}}
\]
where $\hat{P}_{\mathrm{error}}$ is the normalized centralized error
bound, $\rho_{\mathrm{seq}}$ is the effective sequential overhead,
$\delta_{\mathrm{density}}$ is the normalized planning-cost ratio
(tool density), and $M_{\mathrm{topo}} \in \{0,1\}$ is a binary
topology-mismatch flag raised when the detected topology class diverges
from the recommended class.  The weights are $w_1 = 0.4$, $w_2 = 0.3$,
$w_3 = 0.2$, $w_4 = 0.1$.

These bounds are consistent with the scaling laws reported by Kim et
al.~\cite{kim2025scaling}: their empirical observation that error rates
grow with chain depth and that coordination overhead dominates at high
agent counts corresponds to the Error Amplification and Tool Density
bounds respectively.

\subsection{TLA+ Verified Safety Invariants}\label{sec:tla}

Four TLA+ specifications model critical protocols in the convergence
architecture.  Each is verified by the TLC model checker with small-model
parameters (2--3 agents, 2--3 templates/tools, bounded iteration
counts).  Table~\ref{tab:tla} summarizes the specifications and their key
invariants.

\begin{table}[ht]
\centering
\caption{TLA+ specifications and verified safety properties.}
\label{tab:tla}
\small
\begin{tabular}{@{}lll@{}}
\toprule
\textbf{Specification} & \textbf{Key Safety Invariants} & \textbf{Actions} \\
\midrule
TemplateExchangeProtocol
  & Adoption respects \texttt{min\_stage}
  & Export, Import, \\
  & Trust changes only via RecordOutcome
  & RecordOutcome, \\
  & No self-adoption
  & StageAdvance \\
\addlinespace
DevelopmentalGating
  & Critical periods never reopen
  & Tick, AcquireTool, \\
  & Stages monotonically advance
  & OpenPeriod, \\
  & Tools respect capability gates
  & ClosePeriod, Scaffold \\
\addlinespace
ConvergenceDetection
  & Halt is terminal
  & StageResult \\
  & Intervention rate triggers halt
  & (with nondeterministic \\
  &
  & intervention) \\
\addlinespace
EvolutionGating
  & Score monotonically non-decreasing
  & Evolve, Gate, \\
  & Workspace advances only via Gate
  & Rollback \\
  & Version bounded by \texttt{MAX\_VERSIONS}
  & \\
\bottomrule
\end{tabular}
\end{table}

\paragraph{TemplateExchangeProtocol.}
Models cross-agent template sharing with trust-weighted adoption.
An agent imports a peer's template only when three guards are satisfied:
(i) the agent's developmental stage meets the template's minimum stage,
(ii) trust in the peer exceeds \texttt{MIN\_TRUST}, and (iii) the
effective score $\text{peer\_success\_rate} \times \text{trust} \geq
\text{ADOPTION\_THRESHOLD}$.  Trust updates use exponential moving
average (EMA) smoothing, applied exclusively in the \texttt{RecordOutcome}
action---all other actions leave trust unchanged, a structural encoding of
the invariant.

\paragraph{DevelopmentalGating.}
Models lifecycle progression with critical periods and capability gating.
Agents progress through lifecycle stages, advancing through four stages
(Embryonic $\to$ Juvenile $\to$ Adolescent $\to$ Mature).  The key
safety property is \emph{critical period irreversibility}: once a period
transitions from ``open'' to ``closed,'' it can never reopen.  This models time-limited configuration windows that close permanently.

\paragraph{ConvergenceDetection.}
Models the watcher's convergence protocol.  As agents execute stages,
the watcher tracks the ratio of interventions (RETRY, ESCALATE, HALT) to
completed stages.  When this ratio exceeds \texttt{MAX\_RATE}, the
agent is halted in the same atomic step---a safety property expressing
that non-convergent execution is always detected and terminated.

\paragraph{EvolutionGating.}
Models the A-Evolve mutation loop.  The \texttt{Gate} action accepts a
pending mutation only when its benchmark score meets or exceeds the current
score plus \texttt{MIN\_IMPROVEMENT}.  This ensures monotonic score
progression: the agent system never regresses.  The specification intentionally
provides no liveness guarantee, since evolution liveness depends on the
score generator---an external assumption outside the structural model.

\subsection{Co-Design Convergence}\label{sec:codesign}

We apply a simplified version of Zardini's co-design
framework~\cite{zardini2023codesign}, using ordinary function composition
rather than the full categorical apparatus.  We model each adapter as a
\emph{design problem} (DP)---a monotone map from a resource poset to a
functionality poset.  Monotonicity means that more resources yield at
least as many functionalities.

\begin{definition}[Design Problem]
A design problem is a pair $(f, g)$ where $f: \mathcal{R} \to \mathcal{F}$
maps resources to functionalities and $g: \mathcal{R} \to \{0, 1\}$ is a
feasibility predicate.  The map $f$ is monotone: $r_1 \leq r_2 \implies
f(r_1) \leq f(r_2)$.
\end{definition}

\begin{definition}[Series Composition]
Given design problems $\mathsf{DP}_1$ and $\mathsf{DP}_2$, their series
composition $\mathsf{DP}_1 \to \mathsf{DP}_2$ is defined by
$f_{1\to2}(r) = f_2(f_1(r))$, with feasibility
$g_{1\to2}(r) = g_1(r) \wedge g_2(f_1(r))$.
\end{definition}

\begin{definition}[Parallel Composition]
The parallel composition $\mathsf{DP}_1 \| \mathsf{DP}_2$ is defined by
$f_{1\|2}(r) = f_1(r) \cup f_2(r)$, with feasibility
$g_{1\|2}(r) = g_1(r) \wedge g_2(r)$.
\end{definition}

The full adapter stack is a series composition: parsing $\to$ analysis
$\to$ template generation $\to$ scoring.  The adaptive assembly loop
(run $\to$ record $\to$ score $\to$ select) is feedback composition,
where the output of the scoring step feeds back as input to the next
iteration's selection.

\begin{theorem}[Adaptive Loop Approximate Convergence]
Let $\mathsf{DP}$ be the feedback composition of the adapter stack, with
scoring function $s: \mathcal{F} \to [0, 1]$ and feedback map
$\phi(r, f) = r \cup \{s(f)\}$.  If $s$ is monotone, then the sequence
$s_0, s_1, s_2, \ldots$ is monotone nondecreasing and bounded above by~$1$.
For any $\varepsilon > 0$, the implementation terminates when
$|s_{n+1} - s_n| < \varepsilon$ or after a configurable maximum number
of iterations.
\end{theorem}

\begin{proof}
The scoring function $s$ maps functionalities to $[0, 1]$.  The feedback
map $\phi$ is monotone by composition of monotone maps.  Therefore the
sequence $s_0 \leq s_1 \leq s_2 \leq \cdots$ is monotone nondecreasing
and bounded above by~$1$, so it converges to $\sup_n s_n = s^*$ by the
monotone convergence theorem.  The implementation's
\texttt{feedback\_fixed\_point()} has two termination modes:
(1)~when $|s_{n+1} - s_n| < \varepsilon$ (default $0.01$) for the
monitored scoring key, the implementation reports \texttt{converged=True},
indicating the tracked value has stabilized within $\varepsilon$;
(2)~when the iteration count reaches the safety bound (default 100)
without satisfying the $\varepsilon$ criterion, the algorithm halts with
\texttt{converged=False} and returns the last computed state.  The
mathematical guarantee (monotone convergence) ensures the sequence
approaches $s^*$; the implementation provides a practical decision
procedure that reports whether that convergence was achieved within
the allocated budget.
\end{proof}

The implementation in \texttt{convergence/codesign.py} provides
\texttt{compose\_series}, \texttt{compose\_parallel}, and
\texttt{feedback\_fixed\_point} with the exact convergence semantics
described above.

%% ============================================================================
%% 4. IMPLEMENTATION
%% ============================================================================
\section{Implementation}\label{sec:impl}

The convergence architecture is implemented in 14 Python modules under
\texttt{operon\_ai/convergence/}, tested by 216 convergence-specific unit
tests as part of a broader suite of 1{,}404 tests across the entire Operon
codebase.  Table~\ref{tab:adapters} summarizes the five adapters.

\begin{table}[ht]
\centering
\caption{Convergence adapters: source frameworks and key functions.}
\label{tab:adapters}
\small
\begin{tabular}{@{}llll@{}}
\toprule
\textbf{Adapter} & \textbf{Source} & \textbf{Parse Function} & \textbf{Template Function} \\
\midrule
Swarms      & Graph workflows     & \texttt{parse\_swarm\_topology}      & \texttt{swarm\_to\_template} \\
DeerFlow    & LangGraph sessions  & \texttt{parse\_deerflow\_session}    & \texttt{deerflow\_to\_template} \\
AnimaWorks  & Supervisor configs  & \texttt{parse\_animaworks\_org}      & \texttt{animaworks\_to\_template} \\
Ralph       & Hat orchestration   & \texttt{parse\_ralph\_config}        & \texttt{ralph\_to\_template} \\
A-Evolve    & Workspace manifests & \texttt{parse\_aevolve\_workspace}   & \texttt{aevolve\_to\_template} \\
\bottomrule
\end{tabular}
\end{table}

\subsection{Bidirectional DeerFlow Skill Bridge}

DeerFlow's Markdown-based skill system is structurally similar to Operon's
\texttt{PatternTemplate}: both represent reusable workflow patterns with
metadata and staged execution steps.  The module
\texttt{convergence/deerflow\_skills.py} provides bidirectional conversion:
\texttt{skill\_to\_template} parses DeerFlow Markdown skills (extracting
frontmatter, workflow steps, and tool references) into scored
\texttt{PatternTemplate} instances, while \texttt{template\_to\_skill}
exports Operon templates as DeerFlow-compatible Markdown.  This enables
cross-framework template exchange: Operon can import DeerFlow's
community-contributed skill library, and DeerFlow can consume Operon's
structurally validated templates.

\subsection{Hybrid Assembly}

The \texttt{hybrid\_skill\_organism} function implements a two-phase
assembly strategy.  In the first phase, it queries the pattern library for
templates matching the task fingerprint; if a template with a sufficiently
high success score exists, it is used directly.  In the second phase
(fallback), an LLM generator produces a one-shot agent configuration, which is
registered in the library for scoring refinement over subsequent runs.
This library-first, generator-fallback design ensures that the system
improves with experience while maintaining availability for novel tasks.

\subsection{Cognitive Extensions}

Three modules extend the cognitive layer:

\begin{itemize}[leftmargin=*,itemsep=2pt]
  \item \textbf{PrimingView} (in \texttt{patterns/priming.py}): a subclass
    of \texttt{SubstrateView} that adds recent outputs, telemetry,
    experience records, and developmental status as multi-channel context.
  \item \textbf{HeartbeatDaemon} (in \texttt{patterns/heartbeat.py}):
    periodic consolidation of short-term observations into long-term
    memory, inspired by AnimaWorks' heartbeat protocol.
  \item \textbf{AsyncOrganizer} (in \texttt{convergence/async\_thinking.py}):
    Fork/Join execution within individual stages, inspired by the
    AsyncThink~\cite{chi2025asyncthink} observation that intra-stage
    parallel reasoning can reduce latency without increasing token cost.
\end{itemize}

\subsection{Catalog Seeding}

The \texttt{convergence/catalog.py} module seeds the pattern library from
multiple sources: \texttt{seed\_library\_from\_swarms} registers built-in
Swarms workflow patterns; \texttt{seed\_library\_from\_deerflow} imports
DeerFlow skills; \texttt{seed\_library\_from\_ralph} converts Ralph hat
configurations; and \texttt{seed\_library\_from\_acg\_survey} registers
templates from the ACG survey's comparison cards~\cite{yue2026acg},
annotated with graph determination time and plasticity mode metadata.
The combined catalog provides a rich initial library for hybrid assembly.

%% ============================================================================
%% 5. RELATED WORK
%% ============================================================================
\section{Related Work}\label{sec:related}

\paragraph{Agent workflow taxonomy.}
The ACG survey by Yue et al.~\cite{yue2026acg} provides a comprehensive
taxonomy of workflow optimization methods, classifying approaches along
two dimensions: graph determination time (offline, pre-execution,
in-execution) and graph plasticity mode (static, dynamic).  Our work
operates within this taxonomy: the adapter architecture is a pre-execution
analysis layer, while the watcher provides in-execution monitoring.  The
ACG framework treats workflow graphs as optimizable structures; we add
formal verification and convergence guarantees that the survey identifies
as open problems.

\paragraph{Plural intelligence.}
Evans et al.~\cite{evans2026plural} argue that the next advance in AI
will be combinatorial societies of specialized agents rather than
monolithic models, drawing on primate social scaling and cultural ratchet
effects.  Our work provides the formal infrastructure---structured composition,
verified safety invariants, convergence proofs---for the ``plural
intelligence'' vision they articulate.

\paragraph{Scaling laws.}
Kim et al.~\cite{kim2025scaling} derive empirical scaling laws for
multi-agent systems, finding that error rates grow with chain depth and
coordination overhead dominates at high agent counts.  These findings
are consistent with our epistemic theorems: the Error Amplification
Bound predicts the chain-depth dependence; the Tool Density bound
predicts the coordination overhead curve.

\paragraph{Intra-stage reasoning.}
Chi et al.'s AsyncThink~\cite{chi2025asyncthink} demonstrates that
allowing agents to reason asynchronously within a stage---forking
parallel thought branches and joining at a synchronization point---reduces
latency without increasing token cost.  Our \texttt{AsyncOrganizer}
formalizes this pattern with configurable concurrency ratios and
critical-path scheduling.

\paragraph{Co-design theory.}
Zardini~\cite{zardini2023codesign} develops a general theory of co-design
based on monotone maps between resource and functionality posets, with
series, parallel, and feedback composition operators.  We apply this
framework to model each adapter as a design problem and prove that the
adaptive assembly loop converges to a fixed point.

%% ============================================================================
%% 6. DISCUSSION + FUTURE WORK
%% ============================================================================
\section{Discussion and Future Work}\label{sec:discussion}

\paragraph{From analysis to compilation.}
The current adapter architecture analyzes structure but does not generate
executable workflows.  The natural next step is \emph{compilers}: functions
like \texttt{organism\_to\_swarms} and \texttt{organism\_to\_deerflow} that
compile Operon agent systems into Swarms workflow configurations or DeerFlow
LangGraph sessions.  For DeerFlow, the deepest integration is a native
LangGraph watcher node (\texttt{operon\_watcher\_node}) that provides
real-time structural monitoring inside DeerFlow's execution loop.
A Ralph compiler (\texttt{organism\_to\_ralph}) would emit hat
configurations with backpressure constraints derived from the epistemic
analysis.

\paragraph{Cross-target evaluation.}
A systematic evaluation harness running identical tasks across Operon,
Swarms, DeerFlow, AnimaWorks, and Ralph would provide empirical validation
of the structural predictions.  The ACG survey's recommendation to
measure structural variation across inputs is a valuable addition to such
an evaluation protocol, alongside standard metrics (success rate, token
cost, latency, intervention count, convergence rate).

\paragraph{Prompt optimization.}
The adapter architecture optimizes topology but leaves node-level content
(prompts, demonstrations) fixed.  DSPy-style prompt optimization
integrated via the A-Evolve mutation loop could address this gap: the
evolutionary Solve--Observe--Evolve--Gate--Reload cycle naturally
accommodates prompt mutations as a special case of workspace evolution,
with the fitness gate ensuring monotonic improvement.

\paragraph{Limitations.}
The adapters analyze structure, not runtime behavior.  An architecture that
scores well on all four epistemic bounds may still fail due to prompt
quality, model capability, or environmental non-stationarity.  The TLA+
specifications are verified with small-model parameters; scaling to
production-size state spaces would require compositional verification
techniques or abstraction refinement.  The co-design convergence proof
assumes monotonicity of the scoring function, which may not hold under
distribution shift in the task stream.

%% ============================================================================
%% 7. CONCLUSION
%% ============================================================================
\section{Conclusion}\label{sec:conclusion}

We have presented a structured adapter architecture for integrating Operon's
structural analysis layer with five external agent orchestration
systems---Swarms, DeerFlow, AnimaWorks, Ralph, and A-Evolve---through a
single intermediate type, \texttt{ExternalTopology}.  The architecture
contributes four epistemic theorems that serve as a structural linter, four
TLA+-verified safety invariants, and a co-design convergence proof showing
that the adaptive assembly loop stabilizes within $\varepsilon$.  By
operating on serializable dicts with zero external imports, the adapter
architecture achieves framework-agnostic analysis while preserving a
structured interchange format with top-level frozen dataclass fields
(nested dicts remain mutable by Python convention).  The 14 convergence modules, 216 convergence tests, 4 TLA+
specifications, and 98 worked examples provide a concrete, verifiable
foundation for multi-agent system integration.

%% ============================================================================
%% BIBLIOGRAPHY
%% ============================================================================
\newpage
\bibliographystyle{plain}
\bibliography{../references}

\end{document}
